Este anexo presenta el desglose exhaustivo de los requisitos funcionales y no funcionales del proyecto \enquote{Asistente de Preparación \gls{pmp}}. Siguiendo la metodología ágil \gls{scrum} adaptada para este trabajo, el alcance se ha organizado en \textbf{épicas} (grandes contenedores de funcionalidad) e \textbf{historias de usuario} (unidades de trabajo implementables), cada una con sus respectivos \textbf{criterios de aceptación}.

Esta documentación sirvió como el \textit{product backlog} vivo del proyecto, guiando las decisiones de diseño arquitectónico (Capítulo 4) y la implementación técnica (Capítulo 5).

\section{Mapa de épicas}

El sistema se estructura en torno a cinco ejes funcionales principales (ver Tabla \ref{tab:epicas}):

\begin{table}[h]
\centering
\begin{tabularx}{0.95\textwidth}{|l|p{4cm}|X|l|}
\hline
\textbf{ID} & \textbf{Épica} & \textbf{Descripción} & \textbf{Prioridad} \\
\hline
\textbf{E01} & \textbf{Gestión de identidad y personalización} & Funcionalidades para el registro, autenticación, gestión de perfil y personalización de la experiencia del usuario (onboarding). & Alta \\
\hline
\textbf{E02} & \textbf{Núcleo de aprendizaje conversacional (\gls{ia})} & El motor de chat inteligente, incluyendo la gestión de prompts dinámicos, memoria de contexto y los múltiples roles pedagógicos del asistente. & Crítica \\
\hline
\textbf{E03} & \textbf{Simulación y evaluación} & Herramientas para la generación procedural de exámenes, ejecución de simulacros cronometrados y análisis de resultados. & Crítica \\
\hline
\textbf{E04} & \textbf{Gamificación y progresión} & Mecánicas de juego (niveles, rachas, logros) diseñadas para fomentar el hábito de estudio y visualizar el avance. & Media \\
\hline
\textbf{E05} & \textbf{Arquitectura, seguridad y calidad} & Requisitos no funcionales transversales: rendimiento, privacidad de datos, seguridad de \gls{api} y manejo de errores. & Alta \\
\hline
\end{tabularx}
\caption{Resumen de épicas del proyecto}
\label{tab:epicas}
\end{table}

\section{Detalle de historias de usuario}

A continuación, se detallan las historias de usuario que componen cada épica.

\subsection{Épica 1: Gestión de identidad y personalización (E01)}

\subsubsection{HU-01: Onboarding de nuevo usuario (tutorial)}
\textbf{Como} nuevo estudiante,
\textbf{quiero} ser guiado a través de un recorrido interactivo la primera vez que ingreso,
\textbf{para} comprender rápidamente la metodología de \enquote{Mundos} y cómo utilizar las herramientas de \gls{ia} sin necesidad de leer un manual.

\textbf{Criterios de aceptación:}
\begin{itemize}
\item El sistema debe detectar si es el primer inicio de sesión del usuario.
\item Se debe desplegar un modal multipaso (wizard) que explique: 1) estructura del \gls{pmbok}, 2) uso del chat, 3) simulador.
\item El sistema debe solicitar el nombre preferido del usuario para personalizar los saludos de la \gls{ia}.
\item Al finalizar, se debe persistir el estado `hasOnboarded: true` en la base de datos para no repetir el tour en futuras sesiones.
\end{itemize}

\subsubsection{HU-02: Autenticación y persistencia de sesión}
\textbf{Como} estudiante recurrente,
\textbf{quiero} mantener mi sesión activa incluso si recargo la página o cierro el navegador,
\textbf{para} no tener que ingresar mis credenciales cada vez que quiero estudiar unos minutos.

\textbf{Criterios de aceptación:}
\begin{itemize}
\item Implementación de autenticación mediante \gls{jwt}.
\item El token debe almacenarse de forma segura en el almacenamiento local o cookies httpOnly.
\item El middleware de la aplicación debe verificar la validez del token en cada cambio de ruta.
\item Si el token expira, el usuario debe ser redirigido automáticamente a la pantalla de Login.
\end{itemize}

\subsubsection{HU-03: Gestión de perfil de usuario}
\textbf{Como} estudiante,
\textbf{quiero} poder ver y editar mis datos básicos (avatar, nombre),
\textbf{para} sentir que el entorno de estudio es personal y propio.

\textbf{Criterios de aceptación:}
\begin{itemize}
\item Visualización del avatar del usuario en la barra lateral y en los mensajes del chat.
\item Integración con el servicio de avatares de PocketBase o uso de \gls{ui} Avatars por defecto.
\item Opción de \enquote{Cerrar Sesión} claramente visible y accesible.
\end{itemize}

\subsection{Épica 2: núcleo de aprendizaje conversacional (E02)}

\subsubsection{HU-04: Chat con streaming en tiempo real}
\textbf{Como} estudiante,
\textbf{quiero} ver la respuesta de la \gls{ia} generándose progresivamente (efecto máquina de escribir),
\textbf{para} reducir la ansiedad de espera y leer a velocidad natural.

\textbf{Criterios de aceptación:}
\begin{itemize}
\item Implementación de `ReadableStream` en el cliente y servidor.
\item La latencia inicial debe ser inferior a 1 segundo.
\item El chat debe realizar \textit{auto-scroll} inteligente (solo si el usuario está al final de la conversación).
\item Indicador visual de \enquote{La \gls{ia} está escribiendo...} durante la generación.
\end{itemize}

\subsubsection{HU-05: Renderizado de contenido enriquecido (markdown)}
\textbf{Como} estudiante,
\textbf{quiero} que las explicaciones incluyan formato (negritas, listas y tablas),
\textbf{para} facilitar la lectura y comprensión de conceptos estructurados.

\textbf{Criterios de aceptación:}
\begin{itemize}
\item El componente de chat debe parsear markdown estándar.
\item Las tablas deben renderizarse con estilos \gls{css} claros y bordes definidos.
\item Soporte para listas anidadas y citas en bloque.
\end{itemize}

\subsubsection{HU-06: Memoria contextual de la conversación}
\textbf{Como} estudiante,
\textbf{quiero} hacer preguntas de seguimiento (ej. \enquote{¿Me das otro ejemplo de eso?}),
\textbf{para} profundizar en un tema sin tener que repetir el contexto anterior.

\textbf{Criterios de aceptación:}
\begin{itemize}
\item El backend debe recibir una ventana deslizante de los últimos N mensajes (mínimo 10).
\item La \gls{ia} debe ser capaz de resolver referencias anafóricas (\enquote{eso}, \enquote{el anterior}).
\end{itemize}

\subsubsection{HU-07: Modo tutor socrático}
\textbf{Como} estudiante que quiere profundizar,
\textbf{quiero} activar un modo donde la \gls{ia} no me dé respuestas directas, sino que me haga preguntas,
\textbf{para} desarrollar mi propio razonamiento crítico.

\textbf{Criterios de aceptación:}
\begin{itemize}
\item \textit{Prompt de sistema} específico que instruya a la \gls{ia} a responder \textit{siempre} con una pregunta guía.
\item La \gls{ia} debe evaluar la respuesta del usuario y guiarlo hacia la conclusión correcta paso a paso.
\item Activación seleccionable desde el menú de herramientas.
\end{itemize}

\subsubsection{HU-08: Modo roleplay (simulación de crisis)}
\textbf{Como} futuro gerente de proyectos,
\textbf{quiero} practicar mis habilidades blandas interactuando con \enquote{stakeholders virtuales},
\textbf{para} aprender a manejar conflictos y negociaciones difíciles.

\textbf{Criterios de aceptación:}
\begin{itemize}
\item La \gls{ia} debe adoptar una \enquote{persona} (ej. cliente enojado, patrocinador impaciente).
\item El escenario debe presentar un conflicto realista de proyecto.
\item La \gls{ia} debe reaccionar emocionalmente a las respuestas del usuario (calmarse si se gestiona bien, escalar si no).
\end{itemize}

\subsubsection{HU-09: Modo taller de entregables (workshop)}
\textbf{Como} estudiante práctico,
\textbf{quiero} ayuda para redactar documentos reales (ej. el acta de constitución del proyecto),
\textbf{para} entender la estructura y contenido de los artefactos del \gls{pmbok}.

\textbf{Criterios de aceptación:}
\begin{itemize}
\item La \gls{ia} debe guiar la redacción sección por sección.
\item Debe ofrecer ejemplos de redacción para cada apartado.
\item Al final, debe compilar el documento completo en formato Markdown.
\end{itemize}

\subsubsection{HU-10: Modo \enquote{abogado del diablo} (debate)}
\textbf{Como} estudiante avanzado,
\textbf{quiero} debatir contra la \gls{ia} sobre temas polémicos de gestión,
\textbf{para} reforzar mis argumentos y convicciones éticas.

\textbf{Criterios de aceptación:}
\begin{itemize}
\item La \gls{ia} debe tomar deliberadamente una postura contraria o cuestionable (pero plausible).
\item Debe desafiar los argumentos del usuario basándose en falacias comunes o presiones del mundo real.
\item Debe conceder la victoria si el usuario argumenta sólidamente usando el \gls{pmbok}/código de ética.
\end{itemize}

\subsection{Épica 3: simulación y evaluación (E03)}

\subsubsection{HU-11: Generación procedural de exámenes}
\textbf{Como} estudiante,
\textbf{quiero} generar exámenes de práctica ilimitados sobre temas específicos,
\textbf{para} no depender de un banco de preguntas finito que termine memorizando.

\textbf{Criterios de aceptación:}
\begin{itemize}
\item Endpoint de \gls{api} que acepte parámetros: `topic`, `questionCount`, `difficulty`.
\item Prompt de \gls{ia} diseñado para generar salida estrictamente en formato \gls{json}.
\item Las preguntas generadas deben ser inéditas (creadas en el momento).
\item Validación de estructura \gls{json} antes de presentar al usuario.
\end{itemize}

\subsubsection{HU-12: Simulador de examen (interfaz de examen)}
\textbf{Como} estudiante,
\textbf{quiero} una interfaz libre de distracciones con un temporizador,
\textbf{para} simular las condiciones de presión del examen real.

\textbf{Criterios de aceptación:}
\begin{itemize}
\item Ocultar barras de navegación y chat durante el examen.
\item Temporizador decreciente visible.
\item Navegación secuencial (anterior/siguiente) y mapa de preguntas (acceso directo).
\end{itemize}

\subsubsection{HU-13: Evaluación y feedback detallado}
\textbf{Como} estudiante,
\textbf{quiero} recibir una calificación inmediata y explicaciones de cada pregunta al finalizar,
\textbf{para} aprender de mis errores al instante.

\textbf{Criterios de aceptación:}
\begin{itemize}
\item Cálculo automático de puntaje.
\item Desglose de respuestas correctas e incorrectas.
\item Explicación justificada para la opción correcta.
\item Explicación de por qué cada distractor (opción incorrecta) es erróneo.
\end{itemize}

\subsubsection{HU-14: Historial de simulaciones}
\textbf{Como} estudiante,
\textbf{quiero} acceder a un registro de mis exámenes pasados,
\textbf{para} revisar mi evolución y repasar preguntas falladas anteriormente.

\textbf{Criterios de aceptación:}
\begin{itemize}
\item Listado de exámenes en el dashboard con fecha y puntaje.
\item Posibilidad de abrir un examen pasado en modo \enquote{solo lectura} para revisión.
\item Persistencia del objeto \gls{json} completo del examen (para no perder las preguntas generadas).
\end{itemize}

\subsection{Épica 4: gamificación y progresión (E04)}

\subsubsection{HU-15: Sistema de bloqueo de niveles (ruta de aprendizaje)}
\textbf{Como} estudiante novato,
\textbf{quiero} que el contenido se desbloquee progresivamente,
\textbf{para} seguir un orden lógico y no abrumarme.

\textbf{Criterios de aceptación:}
\begin{itemize}
\item Visualización de niveles como nodos en un mapa o lista.
\item Estado visual claro: \enquote{completado} (verde), \enquote{disponible} (azul), \enquote{bloqueado} (gris/candado).
\item Lógica de desbloqueo: nivel N disponible solo si nivel N-1 completado.
\end{itemize}

\subsubsection{HU-16: Dashboard de métricas de desempeño}
\textbf{Como} estudiante,
\textbf{quiero} ver mi rendimiento desglosado por los 3 dominios del examen (personas, procesos, negocio),
\textbf{para} identificar mis áreas débiles.

\textbf{Criterios de aceptación:}
\begin{itemize}
\item Gráficos de barra o radar por dominio.
\item Cálculo basado en el histórico de preguntas respondidas en simulaciones.
\item Indicadores de tendencia (mejorando/empeorando).
\end{itemize}

\subsection{Épica 5: arquitectura, seguridad y calidad (E05)}

\subsubsection{HU-17: Privacidad de datos (multi-inquilino lógico)}
\textbf{Como} usuario consciente de la privacidad,
\textbf{quiero} que mis conversaciones y errores sean absolutamente privados,
\textbf{para} estudiar con confianza y sin miedo al juicio.

\textbf{Criterios de aceptación:}
\begin{itemize}
\item Implementación estricta de \textbf{\gls{rls}} en la base de datos.
\item Reglas de \gls{api} que impidan físicamente leer registros donde \texttt{user\_id != auth.id}.
\item Validación de seguridad en el backend antes de procesar cualquier solicitud.
\end{itemize}

\subsubsection{HU-18: Protección de infraestructura (\gls{api} Keys)}
\textbf{Como} administrador,
\textbf{quiero} que la llave de \gls{api} de Google Gemini esté oculta y protegida,
\textbf{para} evitar robos de cuota y costos inesperados.

\textbf{Criterios de aceptación:}
\begin{itemize}
\item Uso exclusivo de \texttt{GOOGLE\_API\_KEY} en el entorno de servidor (Node.js).
\item Prohibición de exponer la variable con prefijo \texttt{NEXT\_PUBLIC\_}.
\item Proxy de \gls{api} (\texttt{/api/chat}) que actúa como intermediario único.
\end{itemize}

\subsubsection{HU-19: Manejo de errores y resiliencia}
\textbf{Como} usuario,
\textbf{quiero} que el sistema me informe amigablemente si algo falla (ej. \gls{ia} no responde),
\textbf{para} no quedarme esperando indefinidamente ante una pantalla congelada.

\textbf{Criterios de aceptación:}
\begin{itemize}
\item Bloques \texttt{try-catch} en todas las llamadas asíncronas.
\item Mensajes de error legibles por humanos (\enquote{La IA está tardando demasiado}, \enquote{Error de conexión}).
\item Capacidad de reintentar una acción fallida sin recargar toda la página.
\end{itemize}

\subsubsection{HU-20: Diseño responsivo (mobile first)}
\textbf{Como} estudiante ocupado,
\textbf{quiero} poder estudiar desde mi celular en el transporte público,
\textbf{para} aprovechar los tiempos muertos.

\textbf{Criterios de aceptación:}
\begin{itemize}
\item Interfaz adaptable a pantallas móviles (sidebar colapsable).
\item Tamaños de fuente y botones táctiles adecuados (min 44px).
\item El chat y el simulador deben ser plenamente funcionales en viewport móvil.
\end{itemize}

\subsubsection{HU-21: Validación de integridad de datos}
\textbf{Como} sistema,
\textbf{quiero} asegurar que los datos generados por la \gls{ia} tengan el formato correcto,
\textbf{para} evitar que la aplicación se rompa al intentar mostrar una pregunta mal formada.

\textbf{Criterios de aceptación:}
\begin{itemize}
\item Parsing robusto de \gls{json} proveniente de la \gls{ia} (manejo de posibles caracteres extraños).
\item Verificación de existencia de campos obligatorios (\texttt{question}, \texttt{options}, \texttt{answer}) antes de renderizar.
\item Descartar pregunta o regenerar si la validación falla.
\end{itemize}
